% Two-level translation: 32-bit virtual address split into a 10-bit outer
% index p1, a 10-bit inner index p2 and a 12-bit offset.  Outer entry p1
% points to an inner table, whose entry p2 holds the PFN.
\begin{tikzpicture}[x=1cm,y=1cm, font=\footnotesize\sffamily, >={Stealth[length=1.8mm]},
  bits/.style={font=\scriptsize\sffamily, text=ln-muted, inner sep=1pt},
  cellt/.style={font=\scriptsize\ttfamily},
  ttl/.style={font=\sffamily\scriptsize\bfseries, anchor=south, inner sep=1pt}]
  % --- virtual address
  \draw[fill=ln-accent!15] (0,0) rectangle (2.4,0.55);
  \node at (1.2,0.275) {$p_1 = \texttt{1}$};
  \draw[fill=ln-box] (2.4,0) rectangle (4.8,0.55);
  \node at (3.6,0.275) {$p_2 = \texttt{0x00A}$};
  \draw[fill=white] (4.8,0) rectangle (7.2,0.55);
  \node at (6.0,0.275) {\iflangja{オフセット}{offset}\enspace\texttt{0x7C6}};
  \node[bits, anchor=south] at (1.2,0.55) {31\,\dots\,22};
  \node[bits, anchor=south] at (3.6,0.55) {21\,\dots\,12};
  \node[bits, anchor=south] at (6.0,0.55) {11\,\dots\,0};
  \node[anchor=west] at (7.3,0.275) {\iflangja{仮想アドレス}{virtual address}};
  % --- outer table: V | inner table
  \node[ttl] at (1.7,-0.92) {\iflangja{外側テーブル}{outer table}};
  \node[bits, anchor=south] at (1.1,-1.2) {V};
  \node[bits, anchor=south] at (1.9,-1.2) {\iflangja{内側}{inner}};
  \fill[ln-accent!15] (0.9,-1.62) rectangle (2.5,-2.04);
  \foreach \r/\v/\p/\lab in {0/1/A/0, 1/1/B/{}, 2/1/C/2, 3/0/{{\sffamily ---}}/3, 4/{}/{}/{}, 5/1/D/1023} {
    \draw (0.9,-1.2-0.42*\r) rectangle (1.3,-1.62-0.42*\r);
    \draw (1.3,-1.2-0.42*\r) rectangle (2.5,-1.62-0.42*\r);
    \node[cellt] at (1.1,-1.41-0.42*\r) {\v};
    \node[cellt] at (1.8,-1.41-0.42*\r) {\p};
    \node[bits, anchor=east] at (0.85,-1.41-0.42*\r) {\lab};
  }
  \node[font=\small, text=ln-muted, inner sep=0pt] at (1.1,-1.41-0.42*4) {$\vdots$};
  \node[font=\small, text=ln-muted, inner sep=0pt] at (1.9,-1.41-0.42*4) {$\vdots$};
  \draw[->] (0.15,0) -- (0.15,-1.83) -- (0.9,-1.83);
  % index of outer entry 1, above the p1 arrow that enters its row
  \node[bits, anchor=south east] at (0.85,-1.78) {1};
  % --- inner table B: V | PFN, top edge at the height of outer entry 1
  \node[ttl] at (4.35,-1.55) {\iflangja{内側テーブル B}{inner table B}};
  \node[bits, anchor=south] at (3.65,-1.83) {V};
  \node[bits, anchor=south] at (4.6,-1.83) {PFN};
  \fill[ln-accent!15] (3.4,-3.09) rectangle (5.3,-3.51);
  \foreach \r/\v/\pfn/\lab in {0/1/0x2A017/0, 1/0/{{\sffamily ---}}/1, 2/{}/{}/{}, 3/1/0x1F3C5/10, 4/{}/{}/{}, 5/1/0x0C4E2/1023} {
    \draw (3.4,-1.83-0.42*\r) rectangle (3.9,-2.25-0.42*\r);
    \draw (3.9,-1.83-0.42*\r) rectangle (5.3,-2.25-0.42*\r);
    \node[cellt] at (3.65,-2.04-0.42*\r) {\v};
    \node[cellt] at (4.6,-2.04-0.42*\r) {\pfn};
    \node[bits, anchor=east] at (3.35,-2.04-0.42*\r) {\lab};
  }
  \foreach \r in {2,4} {
    \node[font=\small, text=ln-muted, inner sep=0pt] at (3.65,-2.04-0.42*\r) {$\vdots$};
    \node[font=\small, text=ln-muted, inner sep=0pt] at (4.6,-2.04-0.42*\r) {$\vdots$};
  }
  % --- pointer from outer entry 1 to the start of inner table B
  \fill (2.3,-1.83) circle (1pt);
  \draw[->] (2.3,-1.83) -- (3.4,-1.83);
  % --- p2 selects an entry of B
  \draw[->] (4.4,0) -- (4.4,-0.35) -- (5.75,-0.35) -- (5.75,-3.30) -- (5.3,-3.30);
  % --- physical address
  \draw[->] (4.6,-4.35) -- (4.6,-5.0);
  \draw[fill=ln-accent!15] (2.4,-5.55) rectangle (4.8,-5.0);
  \node at (3.6,-5.275) {PFN\enspace\texttt{0x1F3C5}};
  \draw[fill=white] (4.8,-5.55) rectangle (7.2,-5.0);
  \node at (6.0,-5.275) {\iflangja{オフセット}{offset}\enspace\texttt{0x7C6}};
  \node[bits, anchor=north] at (3.6,-5.55) {29\,\dots\,12};
  \node[bits, anchor=north] at (6.0,-5.55) {11\,\dots\,0};
  \node[anchor=west] at (7.3,-5.275) {\iflangja{物理アドレス}{physical address}};
  % --- offset is copied
  \draw[->] (6.8,0) -- (6.8,-5.0);
\end{tikzpicture}
